\section{Problem Formulation}
\label{sec:formulation}

The Vehicle Routing Problem with Time Windows (VRPTW) is formally defined on a complete directed graph $\mathcal{G} = (\mathcal{V}, \mathcal{A})$. The vertex set $\mathcal{V} = \mathcal{N} \cup \{0, n+1\}$ comprises customer nodes $\mathcal{N} = \{1, 2, \dots, n\}$, origin depot node $0$, and destination depot node $n+1$. The physical depot is modeled as two separate nodes ($0$ and $n+1$) to represent route departure and return without self-loops. The arc set is $\mathcal{A} = \{(i, j) : i, j \in \mathcal{V}, i \neq j, i \neq n+1, j \neq 0\}$.

Each customer $i \in \mathcal{N}$ has a non-negative demand $q_i > 0$, service duration $s_i \ge 0$, and service time window $[e_i, l_i]$, where $e_i$ and $l_i$ denote the earliest and latest allowable service start times, respectively. Depot nodes satisfy $q_0 = q_{n+1} = 0$, $s_0 = s_{n+1} = 0$, and $[e_0, l_{n+1}]$ defines the global planning horizon. Each arc $(i, j) \in \mathcal{A}$ has travel distance $c_{ij} \ge 0$ and travel time $t_{ij} \ge 0$. A homogeneous fleet of $K$ vehicles, denoted by $\mathcal{K} = \{1, \dots, K\}$, is stationed at the depot, with each vehicle possessing carrying capacity $Q$.

\subsection{Decision Variables}
The optimization model employs two sets of decision variables:
\begin{itemize}
    \item $x_{ij}^k \in \{0, 1\}$: Binary variable equal to $1$ if vehicle $k \in \mathcal{K}$ traverses arc $(i, j) \in \mathcal{A}$, and $0$ otherwise.
    \item $w_i^k \ge 0$: Continuous variable specifying the time at which vehicle $k \in \mathcal{K}$ begins service at node $i \in \mathcal{V}$.
\end{itemize}

\subsection{Mathematical Model}
The primary objective is lexicographic: first minimizing the number of deployed vehicles ($\NV$), and secondarily minimizing total travel distance ($\TD$):
\begin{equation}
  \operatorname{lex}\,\min \left( \sum_{k \in \mathcal{K}} \sum_{j \in \mathcal{N}} x_{0j}^k, \quad \sum_{k \in \mathcal{K}} \sum_{(i,j) \in \mathcal{A}} c_{ij}\,x_{ij}^k \right)
  \label{eq:lexicographic_obj}
\end{equation}
In computational solver implementations, this hierarchy is evaluated via a scalarized surrogate objective using a large penalty $\mathcal{M} \gg \sum_{(i,j) \in \mathcal{A}} c_{ij}$:
\begin{equation}
  \min \sum_{k \in \mathcal{K}} \left( \mathcal{M} \sum_{j \in \mathcal{N}} x_{0j}^k + \sum_{(i,j) \in \mathcal{A}} c_{ij}\,x_{ij}^k \right)
  \label{eq:scalarized_obj}
\end{equation}
subject to:
\begin{align}
  &\sum_{k \in \mathcal{K}} \sum_{j \in \mathcal{V}, j \neq i} x_{ij}^k = 1, &&\forall i \in \mathcal{N} \label{eq:visit_once} \\
  &\sum_{j \in \mathcal{N} \cup \{n+1\}} x_{0j}^k \le 1, &&\forall k \in \mathcal{K} \label{eq:depot_outflow} \\
  &\sum_{j \in \mathcal{N} \cup \{0\}} x_{j, n+1}^k - \sum_{j \in \mathcal{N} \cup \{n+1\}} x_{0j}^k = 0, &&\forall k \in \mathcal{K} \label{eq:depot_inflow} \\
  &\sum_{i \in \mathcal{V}, i \neq h} x_{ih}^k - \sum_{j \in \mathcal{V}, j \neq h} x_{hj}^k = 0, &&\forall h \in \mathcal{N}, \; k \in \mathcal{K} \label{eq:flow_conservation} \\
  &\sum_{i \in \mathcal{N}} q_i \sum_{j \in \mathcal{V}, j \neq i} x_{ij}^k \le Q, &&\forall k \in \mathcal{K} \label{eq:capacity_constraint} \\
  &w_i^k + s_i + t_{ij} - M_{ij} (1 - x_{ij}^k) \le w_j^k, &&\forall (i, j) \in \mathcal{A}, \; k \in \mathcal{K} \label{eq:time_propagation} \\
  &e_i \le w_i^k \le l_i, &&\forall i \in \mathcal{V}, \; k \in \mathcal{K} \label{eq:time_window_bound} \\
  &x_{ij}^k \in \{0, 1\}, \quad w_i^k \ge 0, &&\forall (i, j) \in \mathcal{A}, \; k \in \mathcal{K} \label{eq:variable_domains}
\end{align}

Constraint~\eqref{eq:visit_once} ensures that every customer is visited exactly once. Constraints~\eqref{eq:depot_outflow} and \eqref{eq:depot_inflow} enforce that each vehicle departs from origin depot $0$ and arrives at destination depot $n+1$ at most once. Flow conservation across all customer nodes is guaranteed by \eqref{eq:flow_conservation}. Constraint~\eqref{eq:capacity_constraint} prevents vehicle capacity violations. Schedule propagation is maintained via \eqref{eq:time_propagation}, where $M_{ij} = \max\{0, l_i + s_i + t_{ij} - e_j\}$ defines the tightest valid Big-$M$ upper bound. Customer time-window compliance is governed by \eqref{eq:time_window_bound}; if a vehicle arrives at customer $i$ prior to $e_i$, it incurs an idle waiting time such that service begins at $w_i^k = e_i$. Finally, \eqref{eq:variable_domains} defines the variable domains.
